\section{Far-source estimates and uniform coupling}
\label{fs:section}

We verify all source and current-correction hypotheses of
\Cref{fe:energy-reduction} for the lab equation of
Lemma~\ref{far:structure} and its differentiated $D$ equation, on an
arbitrary smooth finite slab.  In particular the background coefficients
have the stated symbol and $D$ gains, the exact lab gauge gives the
differentiated weak-null polarization, and
Lemma~\ref{far:boundary-input} supplies the initial and overlap bounds.
The quantities from \Cref{fe:section}, including the news estimate
\eqref{fe:F11}, are consequences of the same small bootstrap used here.
They are not assumptions about an already global evolution.
We write $N$ for a fixed multiple of the total bootstrap size, allowing
also the input size in this coefficient.  The three target levels are
$(U_i,N_i)=(h,9),(h,8),(\psi,7)$ for $i=0,H,P$ respectively.
The coefficients, source splittings and current corrections below are
always expressed in Cartesian laboratory components before conversion
to the emitter frame.

The source budgets are \eqref{fe:F14} and the mixed charge
\eqref{fs:mixed-charge}, with the bulk measure
$\|\cdot\|_{\mathrm b}=L^2(dt\,dr\,d\omega/r)$ from
\Cref{fe:currents}.  For each equation, $\mathcal P$ is the sum of its
corrections with their actual signs.  The mixed charge uses only the
corrected $\xi_X$ cap in \eqref{fe:F12}; it requires no $M$-weighted
bound on an arbitrary uncorrected good derivative.

\subsection{Ordered derivatives and weight comparisons}
For a target level $i$ and $|I|+j\le N_i$, write
$\chi^i_{Ij}=Z^I(rT^jU_i)$ for the corresponding unknown of
\Cref{fe:coordinates}.  Positive sums over spatial-emitter words up to
length $b$, with \emph{fixed} lab-time count $j$, will be used; define
lists, or sizes of lists, by
\begin{equation}
\label{fs:lists}
\begin{aligned}
 V^U_{b j}&=|Z^{\le b}T^j U|,
 &G^U_{b j}&=V^U_{b+1,j},\\
 B^U_{b j}&=G^U_{b j}+r V^U_{b,j+1},
 &&\qquad U=h,\psi .
\end{aligned}
\end{equation}
A superscript can equally be the level used. These letters with derivative indices do not denote the target or its velocity. In bulk \(W_i^{1/2}G^i_{b j},\ w_{i,j}^{1/2}B^i_{b j}\) are controlled by \eqref{fe:F1} with spatial counts \(\le b\). In uses of the time-slice cap, \(B^i_{b j}\) is bounded by \(C\sum_{|I|\le b}|B_1\chi_{I j}|\) plus \(C V^i_{b j}\); we will charge the latter separately there. Similar statements hold for \(G\) with \(G_1\). These facts follow by commuting \(Z\) at fixed \(\partial_q\), using the first derivative formulas. Point entries on spheres, in \(L^2_\omega\), include
\begin{equation}
\label{fs:G1}
\begin{array}{ll}
 V^h_{b j}\lesssim N r^{-1+2\eta_*}\tau^{-A_H/2}&(b+j\le8),\\
 V^h_{b j}\lesssim N r^{-.48}&(b+j\le9),\\
 B^h_{b j},\ r V^\psi_{b j}\lesssim N w_H^{-1/2}&(b+j\le7),\\
 B^\psi_{b j}\lesssim N w_H^{-1/2}&(b+j\le6).
\end{array}
\end{equation}
Use \eqref{fe:F7}, \eqref{fe:F8} (including expanding one more lab \(T\) before the \(Z\)'s for the last bound). Two additional rotations give sup versions. In particular low \(G^h\) has the bound of the first line, and low bad entries of either field have size \(\lesssim N\sqrt{1+M}\).

Some comparisons used below, for large \(r\), are
\begin{equation}
\label{fs:G2}
\begin{aligned}
 \frac{W_i}{w_{k,j}}&\lesssim r^{p_i+.021}\tau^{A_i-A_k},
 &w_{i,j}&\lesssim W_i,\\
 \frac{w_{P,j}}{w_H}&\lesssim l^{A-A_H},
 &\frac{w_H}{w_{P,j}}&\lesssim l^{-(A-A_H)+7d_*},\\
 \frac{w_{0,j}}{w_H}&\lesssim r^{-.018}\mu^{-A_H}.
\end{aligned}
\end{equation}
Write \(m_{\rm cone}\) for the weight \(m=(r^{-1}+l^{-1})^{-1}\) of \Cref{fe:section}, to distinguish it from derivative counts denoted by \(m\) below. For \(i=0,P\), \(\sqrt{w_{i,j}/w_{i,j'}}\lesssim m_{\rm cone}^{9d_*/2}\), and for \(j'=j+1\) the bound is \(C m_{\rm cone}^{-d_*/2}\). Indeed \(q\ge-r\), \(m_-m_+\asymp m_{\rm cone}\). On \(|q|\le C r^\nu\) weights within each level are equivalent in \(j\), and \(w_{P,j}\asymp w_H\). Outside \(|q|\le r^\nu\),
\[
 w_H^{-1/2}\lesssim r^{(e-(\nu-\delta)A_H)/2}l^{-A_H/2}.
\]
We use without further comment \(r\widetilde a\lesssim N\tau^{-A/2}\), \(|a(q)|\lesssim N\langle q\rangle^{-A/2}\) as coefficient bounds. When \(|q|\gtrsim r^\theta\), \(\theta=.45\), both \(m_{\rm cone}^{-d_*/2}\) and \(l^{(-(A-A_H)+7d_*)/2}\) gain more than \(r^{-3\eta_*}\).

Two useful bad-commutation comparisons, with \(c_0\) any bounded smooth coefficient of \(n,s(q)\), are
\begin{equation}
\label{fs:G3}
\begin{split}
 |\partial_q(c_0\partial_q\chi^i_{I j})| \lesssim&
 B^i_{b,j+1}+r^{-1}B^i_{b+1,j}+|a(q)|B^i_{b j},\qquad b=|I|,\\
 B^h_{b,j+1}\lesssim&
 B^\psi_{b j}+r^{-1}
             \sum_{i'\le j,\ b'+i'\le b+j+1} B^h_{b' i'} .
 \end{split}
\end{equation}
Indeed \(\partial_q(rT^j U)=krT^{j+1}U+s^d b_d Z T^j U\) and \(Z\) commutes with \(\partial_q\). For the second line express
\begin{equation}
\label{fs:G4}
\begin{aligned}
T^{j+1} h=(1-s_s\cdot n)^{-1}
    \Bigl(&T^j\psi-r^{-1}s_s^d b_d Z T^j h\\
           &-\sum_{i'<j}c_{i'j}\,(T^{j-i'}s_s^d)
                                  \partial_d T^{i'}h\Bigr),
\end{aligned}
\end{equation}
and likewise for \(T^{j+2}h\), using the smoothing bounds. For example \(Z^{\le b}s_s^d b_d ZT^{j+1}h\) from the latter identity after multiplication by \(r\) costs \(C r^{-1}B^h_{b+1,j}\) (here derivatives distribute on the expression). Extra \(Z\)'s on all bounded coefficients cost at most constants. Constants \(c_{i'j}\) are combinatorial.

\subsection{The complete source catalogue}
\begin{proposition}[Complete far-source estimate]
\label{fs:source-estimate}
For each commuted equation on such a finite slab, the sum of its
budgets \eqref{fe:F14}, together with its mixed charges
\eqref{fs:mixed-charge}, is bounded, for some fixed $c_1>0$, by
\begin{equation}
\label{fs:F15}
 o_{R_0\to\infty}(1)(\mathsf F+\mathsf S_a)
 +C\sum_{\mathrm{prior}}\mathsf F_{\mathrm{prior}}
 +C(R_0)\mathsf K_R
 +C_{\mathrm{fix}}(e_{\mathrm{in}}+\mathsf N^{1+c_1}).
\end{equation}
The finite prior order is: level $P$ before the two $h$ hierarchies;
within each level lower total order first, then more laboratory times
at equal total order, then fewer dilations $S$.  Equal-index norms are
aggregated.  The constants on prior far quantities are independent of
both matching radii.  Small radial gains and nonlinear bootstrap
products may cross this order.  After fixed finite hierarchical
weighting, \Cref{fe:energy-reduction} therefore gives, for some $c_2>0$,
\begin{equation}
\label{fs:F}
 \mathsf F\le o_{R_0\to\infty}(1)(\mathsf F+\mathsf S_a)
 +C(R_0)\mathsf K_R
 +C_{\mathrm{fix}}(e_{\mathrm{in}}+\mathsf N^{1+c_2}).
\end{equation}
The initial and starting-annulus correction costs satisfy the same
bound.  In particular the pure acceleration defect on the starting
annulus has its direct small radial coefficient on $\mathsf S_a$.
\end{proposition}

The following ledger identifies the origin and destination of every term.
The numbered claims retain the coefficient placements, corrections, and
endpoint estimates needed for each row.
\begin{center}
\small
\begin{tabular}{@{}>{\raggedright\arraybackslash}p{.20\linewidth}
                  >{\raggedright\arraybackslash}p{.21\linewidth}
                  >{\raggedright\arraybackslash}p{.24\linewidth}
                  >{\raggedright\arraybackslash}p{.25\linewidth}@{}}
\toprule
Source & Derivative placement & Outgoing treatment & Raw and boundary checks \\
\midrule
Principal commutators & Coefficient counts $\le5$, $\ge6$;
exceptional count $9$ & Currents, direct gain, or mixed caps according
to the placement & \eqref{fs:G3}--\eqref{fs:G6}; Claim 1 \\[3pt]
Additional $D$ commutators & Seven target counts; $B^h$ through eight
& Projected gain and acceleration currents
& \eqref{fs:G7}; Claim 2 \\[3pt]
Pure background defect & Symbol derivatives at every base level
& Direct averaging gain, also on the starting annulus
& \eqref{fs:defect-shell-margins}; Claim 3 \\[3pt]
Linear and suppressed products & Nine counts for $h$; eight $h$ counts
or seven mixed counts for the strong equation
& Radial suppression in both bulk weights
& \eqref{fs:G9}; Claim 3 \\[3pt]
Leading screen quadratic & $h$ levels through $N_i$; strong correction
through eight $B^h$ counts
& Mixed cap for $h$; current \eqref{fs:G12} for $\psi$
& \eqref{fs:G10}--\eqref{fs:G12}; Claims 4--5 \\[3pt]
Correction boundaries & Same time weights; strong products use $2A_H>A$
& Sum signed corrections before forming $\xi$
& Initial symbols and off-trap overlap; Claim 6 \\
\bottomrule
\end{tabular}
\end{center}

\begin{proof}
We verify the catalogue in six claims.  Coefficients denoted by $N$
are bootstrap/input factors, so $N\mathsf F$ is nonlinear.  The pure
background symbol bound $B-\eta=O(r^{-1})$ has no such smallness factor.
All wave equations and laboratory derivatives are taken before the
emitter decomposition.

\par\medskip\noindent\textbf{Claim 1: common principal commutators.}\label{fs:claim-principal}\quad
We first derive the common source forms, then estimate their raw and
outgoing budgets separately.  The distinction is necessary because the
outgoing correction avoids one bad derivative, whereas the raw budget
uses the derivative hierarchy directly.

\smallskip\noindent\emph{Source forms and sphere placements.}
After the flat recurrence in Claim 3 of the proof of
Proposition~\ref{fe:energy-reduction}, every such term is as in \eqref{fe:F5} applied at lower indices \(\ell=(b',j')\), with coefficients
\[
  C_h=Z^d (T^m H)_{\ell_{\rm null}\ell_{\rm null}},
                 \qquad J_h=Z^{\le d} T^m H .
\]
Here \(\ell_{\rm null}=(1,-n)\) (the same \(\ell\) as in \eqref{fe:F5}); the contraction is inside the \(Z\)'s. A length written as an exponent allows any word of that length. The indices, for a target total order \(n_*\) and time count \(j\), have
\[
 j'=j-m,\qquad n'+b_0\le n_*,\quad n'<n_*,
        \qquad n'=b'+j',\quad b_0=d+m.
\]
More explicitly the bad term has the form
\begin{equation}
\label{fs:G5}
C_h\partial_q(c_0\partial_q\chi_\ell)
\end{equation}
where \(c_0\) is as in \eqref{fs:G3}, and spatial counts on \(\chi_\ell\) can be smaller. The other terms, multiplied by \(r\), cost
\begin{equation}
\label{fs:G6}
C |J_h|(B^i_{b'+1,j'}+G^i_{b'+1,j'})
                + C r |a(q) C_h|\, G^i_{b' j'}.
\end{equation}
These are bounds term by term, allowing sums and bounded coefficients. Indeed lab commutation first gives \(P[T^m H]\) on \(rT^{j'}U\) if \(m\ge1\). In the \(Z\) commutations use \(P_g=P_\eta+P[H]\); differences from the flat recurrence differentiate \eqref{fe:F5} by Leibniz (including a scalar multiple from each dilation recurrence), with highest second derivatives on the field cancelling. In particular there is no differentiation of the projector \emph{outside} the full differentiated contraction in \eqref{fs:G5}; differentiations of \(k^{-1}\) by \(Z\) stay inside \(\partial_q\). The mixed and spatial second derivatives and all lower terms obey \eqref{fs:G6}. The \(a(q)\) terms there only differentiate \(s\)-dependent frame factors, with the double contraction of \eqref{fe:F5}.

Products in these formulas and their estimates below use sphere \(L^2_\omega\) on a coefficient if it is high, moving two extra rotations to a low companion; otherwise we use coefficient sup with spare rotations. For \(J_h,C_h\) one may use the latter when \(b_0\le5\); when \(b_0\ge6\) there is room since \(n'+3\le n_*\). The composition sphere bounds follow by the ordinary chain rule with at most one highest factor not in sup. Write \(L_h\) for the bound of \(r|C_h|\) with the indicated placement. For \(b_0\le8\) (without increasing a high index) it is \(O(1+N r^{2\eta_*})\), tight \(O(1)\) on \(|q|\le 10r^\theta\), and \(O(r^{-1+8\eta_*})\) if \(m\ge1\); these use \eqref{fe:F9}, \eqref{fe:F10}, with smallness taken after any fixed bootstrap constants. For \(b_0=9\) use the sphere bounds of \eqref{fe:F10} exactly. Now \eqref{fs:G6} costs an acceptable amount in both raw bulk weights of \eqref{fe:F14}, even ignoring the \(1+M\) denominator. Indeed \(J_h\) costs \(r^{-1+2\eta_*}\) to order eight or \(r^{-.48}\) at nine (then the target is level 0). Use \eqref{fs:G2}. Also \(L_h |a(q)|\lesssim N\) by the tightness and decay split, except possibly at nine; there use \(G^H_{b'j'}\) with \(W_H\) on the low field, so \((W_0/W_H)^{1/2}\) and \(a(q)\) decay give the same sufficiency by \eqref{fe:F10}.

\smallskip\noindent\emph{The raw time-multiplier budget.}
Estimate \eqref{fs:G5} without splitting, namely \(L_h\) times the first line of \eqref{fs:G3} with target square weight \(w_{i,j}\). For \(m\ge1\) it is small by \eqref{fe:F9}, \eqref{fe:F10} and \eqref{fs:G2}, including the \(j\)-transfer. It remains to treat \(m=0\).
At levels \(P,0\), except \(b_0=9\), the \(B_{b',j+1}\) term costs prior norms: \(L_h m_{\rm cone}^{-d_*/2}\) is bounded by tightness and the outer \(m_{\rm cone}\)-gain. The other two terms are small by \(1/r\) and \(a(q)\). 

At level \(H\) instead use both lines of \eqref{fs:G3}; \(B^\psi_{b'j}\) has total order at most seven, and \(L_h(w_H/w_{P,j})^{1/2}\lesssim1\). Angular placements respect this order, and level \(P\) is prior. 

The exceptional coefficient count \(b_0=9\) occurs only at level 0
and has a low input.  Use the same conversion: \eqref{fs:G2} gives \(L_h(w_{0,j}/w_{P,j})^{1/2}\lesssim1\), since the logarithm of \eqref{fe:F10} is offset at small \(|q|\), and the term \((1+\min(|q|,r))^{1/2}r^{-p_0/2}\) is offset too (use \(\mu^{-A_H/2}l^{(-(A-A_H)+7d_*)/2}\) at large \(|q|\)). Similarly \(L_h|a(q)|(w_{0,j}/w_H)^{1/2}\lesssim N\); the \(r^{-1}\) terms are small directly.

\smallskip\noindent\emph{The outgoing budget with a current correction.}
If $b_0\le5$, or if the target is level $P$, use the correction
\[
                         \mathcal P=C_h c_0\partial_q\chi_\ell
\]
for \eqref{fs:G5}, with its actual sign/factor. Corrections for different terms are always summed. The remainder differentiates \(C_h\) by \(\partial_q\), so the bound after multiplying by \(r\) uses \eqref{fe:F9}, cost \(r^{-1+8\eta_*}\), against \(B_1\) in the field; this is small in \eqref{fe:F14} by \eqref{fs:G2}. For the correction norms \(C_h,ZC_h\) cost at most \(O(r^{-1+2\eta_*})\), and when \(Z\) hits the field \(n'+1\le n_*\). In the term \(r d_g\partial_q\mathcal P\) one has \(r d_g\lesssim r^{2\eta_*}\) and uses \eqref{fs:G3}. Thus all its bulk costs are small too; its time cap uses \(B_1\chi_\ell\) alone and is small. Sphere placements for \(b_0\le5\) take \(\partial_q C_h,Z C_h\) with two extra rotations, allowed in \eqref{fe:F9}.

\smallskip\noindent\emph{The outgoing budget with a high coefficient.}
It remains to treat $b_0\ge6$ at levels $0,H$, without a current
correction.  If $m\ge1$, the same raw bound as before, now using \(W_i/w_{i,j'}\) or \(W_i/w_{i,j'+1}\), works with strict gain by \eqref{fe:F9}, \eqref{fe:F10}. 

If \(m=0\), use \eqref{fs:G3} in full, including the two extra rotations on the low input. The terms with \(r^{-1}\) have strict gain by \eqref{fs:G2}. Thus it suffices to bound \(L_h\) times cap/bulk derivatives \(B_1\) of low \(\psi\) fields, and \(L_h |a(q)|\) times \(B_1\) of low \(h\) fields. Indeed the \(V\) pieces in the cap comparison above can instead use \eqref{fe:F1} bulk with \(W_P,W_H\) respectively, by \eqref{fe:F10} (as for \eqref{fs:G6}, and \(L_h(W_i/W_P)^{1/2}\ll 1\)). All low orders used here fit even with two extra rotations. 

On the timelike region \(t\gg r\), the bad bulk weights of levels \(P,H\) suffice for the \(B_1\) costs with gain: use \(w_{P,j}\gtrsim t^A r^{-.025}\), \(W_i\asymp r^{p_i-A_i}t^{A_i}\), and \(w_H\gtrsim t^{A_H}r^{-e-\delta A_H}\) with \(|a(q)|\) decaying; the extreme \(b_0=9\) only occurs for \(W_0\).

On the remaining region $t\lesssim r$, $r\asymp R_f$, use the mixed entry with
$l_{R_f}$ equal to the bounds in \eqref{fe:F9}, \eqref{fe:F10} for
$L_h(w_{P,j}^{-1/2}+|a(q)|w_H^{-1/2})$, supremized at fixed
$\mathfrak u$ on the shell.  Weights and power thresholds may use
$R_f$ in place of $r$, with fixed comparable constants.

\emph{Counts six through eight.} The squared integral of this kernel is
$O(R_f^{.013\eta_*})$.  Indeed
\[
 w_{P,j}^{-1}\lesssim
 r^e\mu^{-A_H}l^{-(A-A_H)+7d_*},
\]
whose integral is $O(r^{e+\delta A_H+\nu(1-A_H)})$ because
$A-7d_*>1$.  The margin needed for the mixed shell charge is
\begin{equation}\label{fs:principal-shell-margin}
 e+\delta A_H+\nu(1-A_H)
 =.011\eta_*+.09\eta_*^2<.013\eta_*<1-p_H.
\end{equation}
On $|q|\gtrsim r^\theta$, the tail integral has enough additional
power gain to offset $L_h^2$; the decay of $a(q)$ handles its entry
as well.  Supremizing uses only comparability of
$r^\delta+|q|$ and $r^\delta+|\mathfrak u|$, so the small-$q$
acceleration bound may be broadened to a band $O(R_f^\delta)$.
This proves summability for $b_0\le8$.

\emph{Count nine.} Only weight $W_0$ occurs.  The squared extra term
$(1+\min(|q|,r))r^{-p_0}$ in $L_h^2$, including the mixed shell
factor squared, contributes at most
\[
 C R_f^{-1}\int_{|\mathfrak u|\lesssim R_f}
 (R_f^\delta+|\mathfrak u|)
 (w_{P,j}^{-1}+|a(q)|^2w_H^{-1})_{\sup}\,du.
\]
Split the $w_P^{-1}$ part at $|\mathfrak u|=R_f^\nu$.  It is
bounded by the sum of
\[
 R_f^{-1+e+\delta A_H+\nu(2-A_H)},\qquad
 R_f^{1-A+7d_*+e+\delta A_H+\nu(A-A_H-7d_*)}.
\]
Both exponents are less than $-1/4$ for the fixed range of $\eta_*$.
For the acceleration part, use $|a|\le CN$ on the broadened band
$|\mathfrak u|\le C R_f^\delta$ and
$|a(q)|^2\le CN^2|\mathfrak u|^{-A}$ outside it.  Its powers are
bounded by
$R_f^{-1+e+2\delta}$ and $R_f^{-1+e+\delta(2-A)}$.
The integral is therefore $O(R_f^{-c_1})$ for a fixed $c_1>0$;
logarithms and $R_f^{8\eta_*}$ losses in the ninth-order coefficient
leave a strict summable margin.  This accounts for every common
principal commutator in both multipliers.

\par\medskip\noindent\textbf{Claim 2: the additional strong-equation terms.}\label{fs:claim-strong}\quad
Before treating \(f_h,D f_h\), the additional strong equation commutators arise already at order zero:
\begin{equation}
\label{fs:G7}
\begin{aligned}
g^{ij}\partial_{ij}\psi
   &=D f_h-(Dg^{ij})\partial_{ij}h\\
   &\quad+2g^{ij}(\partial_i s_s^d)\partial_{jd}h
             +g^{ij}(\partial_{ij}s_s^d)\partial_d h .
\end{aligned}
\end{equation}
In \eqref{fe:F5} on the second term, \(J\) (full) and \(C\) (double contraction) after ordered differentiations through \(b+m\le7\) satisfy in sphere \(L^2\)
\[
           |J|\lesssim N r^{-1+8\eta_*}\tau^{-A_H/2},\qquad
           |C|\lesssim N r^{-2+8\eta_*}\tau^{-A_H/2}.
\]
Indeed \(D B=O(r^{-1}\widetilde a)\) with symbol bounds (freezing makes this derivative zero), \(D h=\psi\). The improved estimate on the double contraction of \(\psi\) uses \(T^m D T^0 h\) via the same formula as \eqref{fs:G4} but solved for \(T^m\psi\), and the projected estimate on \(T^{m+1}h\) in the proof of \eqref{fe:F9}, i.e. without a stationary background term since this is \(h\) itself; \(\Theta\) there costs \(O(r^{-1}\widetilde a)\). Quadratic inverse errors use \eqref{fs:G1}. The full \(r\)-gain on the projected coefficient therefore applies even through all the commutations of \eqref{fs:G7}. After multiplying by \(r\), the bad, mixed/good and \(a(q)\) terms of \eqref{fe:F5} now cost an acceptable amount directly in both raw weights: use \eqref{fs:G3} on \(h\) through total order eight for \(B\), weighted by \(w_H\); \eqref{fs:G2} works because \(A<2A_H\). For \(a(q)\) times good entries use \(W_H\). If the coefficient is high use the two rotations instead on a low input (\(b+m\ge6\) suffices). Likewise the piece \(2H^{ij}\partial_i s_s^d\) in the third term of \eqref{fs:G7} costs \(O(N r^{-2+2\eta_*}\tau^{-A/2})\) in \emph{all} components through these coefficient orders, hence works directly. The last term of \eqref{fs:G7}, after the \(r^2\) source scaling, costs \(C\widetilde a\, B^h\) with bounded coefficient products and likewise works.

Here the derivative accounting is explicit.  If $c=b+m$ derivatives
hit $Dg$ and $n'$ hit $h$, then $c+n'\le7$.
The two physical derivatives on $h$ require a $B^h$ index of at most
$n'+1\le8$, available at level $H$.  For $c\le5$ two additional
rotations on the coefficient stay within its order-seven estimate.
For $c\ge6$, place it in $L^2_\omega$; the input including the two
transferred rotations has $B$ count at most $7-c+1+2\le4$.
For example, comparison of the $J$ or $rC$ bound with the high
$w_H$-weighted bad norm gives the squared factor
\[
 C N^2 r^{p_P-2+16\eta_*+e}\tau^{A-2A_H}.
\]
The radial exponent is $-.38+16.001\eta_*<0$ and
$A-2A_H=-.64+2\eta_*<0$.  These are strict gains before using
the $1+M$ denominator.

For the smooth \(Y=2\eta^{ij}\partial_i s_s^d\) coefficient of \(\partial_{jd}h\), \(J=Z^bT^m Y\) is \(O(r^{-m}\widetilde a)\); \(\partial_q C\) for its differentiated contraction as before costs an extra \(r^{-1}\), since \(\partial_q,Z\) commute and no derivative of the fixed-\(n\) projector occurs. Here \(n'\) in \eqref{fs:G5}, \eqref{fs:G6} for the input \(h\) can equal \(n_*\le7\). In the time multiplier use \eqref{fs:G3} in full for \eqref{fs:G5}: the \(\psi\) term costs \(C r^{1-m}\widetilde a\, (w_{P,j}/w_{P,j-m})^{1/2}\lesssim C N\) times its weighted bad entry, admissible even at the same order. The \(a(q)\) terms on \(B^h,G^h\) use \((w_{P,j}/w_H)^{1/2}\lesssim l^{(A-A_H)/2}\) and acceleration decay. All \(1/r\)-suppressed terms work directly by \eqref{fs:G2}, as do the ordinary mixed/good terms of \eqref{fs:G6}. In the outgoing multiplier use current correction \(C c_0\partial_q\chi^h_\ell\) as for \eqref{fs:G5}. Its \eqref{fe:F14} costs all use \eqref{fs:G2} and \(w_H\), with up to order eight by \eqref{fs:G3}, \(r d_g\lesssim r^{2\eta_*}\); the bulk source remainder costs \(r|\partial_q C|\lesssim\widetilde a\) times \(B_1\). The ordinary mixed/good terms are small too.

One further correction here deals with the explicit \(a(q)\) term of \eqref{fe:F5} (not the one inside \eqref{fs:G5}). More precisely, with \(d_i Z=(b_i-\ell_i s^d b_d/k)Z\) in the first derivative formula, it comes from
\[
         -Y_{\ell_{\rm null}\ell_{\rm null}}\,
                    \partial_q(s^d/k)\, b_d Z(\chi/r)
\]
in \(P[Y]\chi\), and the corresponding ordered differentiations (lab times on \(Y\) and the input taken before \eqref{fe:F5}). It thus has the form \(C\,(\partial_q c_2)\, D_Z\chi_\ell/r\), with sign included in \(D_Z\), an operator of order at most one in \(Z\) with smooth sphere coefficients independent of \(q\); \(c_2\) smooth bounded in \(n,s(q)\). Use correction \(C c_2 D_Z\chi_\ell/r\). The two remainders, multiplied by \(r\), cost \(C'\widetilde a\) times \(B^h,G^h\) entries through order eight, all suppressed by \eqref{fs:G2}. For the bulk correction budget even after \(Z\) it uses \(G^h\) through order eight with the same coefficient bound; \(r d_g\partial_q\) additionally costs \(O(r^{2\eta_*}\widetilde a)\) against \(G^h\) and \(O(r^{-1+2\eta_*}\widetilde a)\) against \(B^h\) (differentiate the formula, including \(|\partial_q c_2|\lesssim |a(q)|\)). These meet \eqref{fe:F14}. For its time-slice budget use \(G_1\chi_\ell\) with \(W_H\); any undifferentiated \(\chi_\ell/r\) term uses \eqref{fs:G1} directly with \(\widetilde a\), giving summability in \(L^2(dr)\) even with \(W_P^{1/2}\). This proves sufficiency for \eqref{fs:G7} except \(D f_h\).

\par\medskip\noindent\textbf{Claim 3: background, linear, and suppressed products.}\label{fs:claim-suppressed}\quad
We finish with the semilinear sources, i.e. linear combinations (from the flat recurrence, thus with constant scalars out front) of
\begin{equation}
\label{fs:G8}
Z^{c_3}(kr T^j f_h),\qquad Z^{c_3}(kr T^j D f_h),
                   \qquad c_3+j\le n_* .
\end{equation}
All product estimates allow indices \emph{up to} the indicated ones. The estimates for \(r\) times these expressions thus concern \(r^2\)-scaled physical sources. Unless specifying a different treatment, a "both weights" bulk estimate below uses \(w_{i,j}^{1/2}\) and \((W_i/(1+M))^{1/2}\) of \eqref{fe:F14}; for suppressed products it suffices simply to use \(W_i^{1/2}\).

For the pure background source, the $r^2$ scaling leaves
$\widetilde a$ at levels $0,H$ and $r^{-1}\widetilde a$ at level $P$;
the latter gain comes from the $D$ derivative of the defect.
We estimate these terms directly, including on the starting annulus.
On a shell $r\asymp R_f$, an enlarged positive majorant is
\begin{equation}\label{fs:acceleration-average}
 A_{R_f}(t)=\frac{C_1}{R_f}
              \int_{t-C_2R_f}^{t}|a(u)|\,du,\qquad a(u)=0\quad(u<0).
\end{equation}
Since $A>1$, Cauchy--Schwarz gives
$\|a\|_{L^1}\le C_A\mathsf S_a$.
The averaging kernel has $L^2$ norm $O(R_f^{-1/2})$, hence
the first estimate below.  For $t\ge2C_2R_f$, the time weight is
comparable throughout the averaging interval; the bounded $L^1$ norm
of the kernel gives the second:
\begin{equation}\label{fs:average-bounds}
 \|A_{R_f}\|_{L^2_t}\le C R_f^{-1/2}\mathsf S_a,
 \qquad
 \|\langle t\rangle^{A/2}A_{R_f}\|_{L^2(t\ge2C_2R_f)}
       \le C\mathsf S_a.
\end{equation}
The fixed intermediate time band is included in the early estimate.
As $A_i\le A$, the late contribution satisfies
\[
 R_f^{p_i-A_i}
 \|\langle t\rangle^{A_i/2}A_{R_f}\|_{L^2(t\ge C R_f)}^2
 \le C R_f^{p_i-A}\mathsf S_a^2 .
\]
The squared early/late shell exponents are $p_i-1,p_i-A$ for $i=0,H$,
with two further inverse powers for $i=P$:
\begin{equation}\label{fs:defect-shell-margins}
\begin{array}{c|cc}
 i&\text{early}&\text{late}\\ \hline
 0&-.36&-.72\\
 H&-\eta_*/6&-.36-\eta_*/6\\
 P&-1.38&-1.74
\end{array}
\end{equation}
and their dyadic sums tend to zero with $R_0$.
This calculation also applies on $[R_0/2,R_0]$: the pure defect is
estimated directly there, without a growing compact coefficient
on $\mathsf S_a$.

For the terms of linear jet size \(O(r^{-2}\partial h+r^{-3}h)\), after multiplication by \(r^2\) one uses \(r^{-1} B^h_{b j'}\), times bounded smooth coefficients with their distributed differentiations. The \(h\)-jet factors in those coefficients cost at most \(C\) in \(L^2_\omega\) even through order nine; in a product with a high such factor, transfer two rotations to the bulk field for the angular placement. More precisely one needs this transfer only at high coefficient orders (say seven and up); all its counts then fit the same level. Using \(w_{i,j'}^{1/2} B^i_{b j'}\) from \eqref{fe:F1}, \eqref{fs:G2} supplies strict radial gain. The leading stationary coefficient constants have symbol bounds independent of the matching radii. In \(D f_h\) these terms give the same estimate for \(\psi\) in place of \(h\), and terms with \(r^{-1}B^h_{b j'}\) accompanied additionally by \(\widetilde a\) or a factor \(V^\psi\) (with distributed indices). For the latter estimates use level \(H\) on \(h\), \eqref{fs:G1} on the extra factors, and \eqref{fs:G2} with \(A<2A_H\). In particular \(\mu\gtrsim\tau\). The orders here sum to at most seven; a high extra factor can be put in sphere \(L^2\) with its stated decay, using spare rotations on the others. To justify this differentiation description one can take \(D\) before the lab \(T^j\), still in Cartesian lab formulas before frame and scale expansions. Linear jet coefficients before smooth field composition consist of background first jets (and gauge coefficients of first-jet size) at order \(O(r^{-2})\), second or first-product jets at \(O(r^{-3})\). On these \(D\) gains \(O(\widetilde a)\) relative to the indicated bounds: \(D B=O(r^{-1}\widetilde a)\) with symbols by \(D x=0\) and the parameter/formal-frame differentiation of the order \(-1\) residual from flat in lab components. Commuting lab derivatives costs \(\partial s_s=O(\widetilde a)\) with symbol bounds; the compensator and its derivatives have the stated stronger bounds from the reduction. Smooth \(h\)-coefficients instead differentiate by \(D h=\psi\). This applies directly in the Ricci and gauge substitution formulas giving \eqref{fe:F3}; it does not require \(D r\) for the emitter radius to vanish.

Here are reusable quadratic product placements. The last term of \eqref{fe:F3} has vanishing coefficient \(\mathcal H=O(r^{-1}+|h|)\) in front of two first jets. After the scaling and differentiations for \(h\), this costs products bounded by
\begin{equation}
\label{fs:G9}
\begin{gathered}
 |Z^b T^m\mathcal H|\, B^h_{b_1 i_1} B^h_{b_2 i_2},\\
 b+m+b_1+i_1+b_2+i_2\le n_* .
\end{gathered}
\end{equation}
with harmless additional smooth sphere coefficients, sums, and analogous metric compositions. Use \(r^{-1+2\eta_*}\) on the coefficient through order eight, \(r^{-.48}\) at nine, in sphere \(L^2\). This also bounds the corresponding rates in sup with two spare rotations. If \(b+m\le6\) take that sup with the better rate, and otherwise put both \(B\) factors in sphere sup using two more rotations on each. Put a lower-order \(B\) in point control \eqref{fs:G1}, the other in bulk \(w_{i,i_k}^{1/2}\) of its occurrence (here \(i=H,0\) is the level, \(i_k\) the relevant time count). There is enough room, since the lower-order factor needs order at most four before the rotations and a coefficient consuming at least seven leaves both factors low. By \eqref{fs:G2}, \eqref{fs:G1}, \eqref{fs:G9} fits both weights with strict radial slack also in the order-nine coefficient case.

For strong estimates a similar sufficient suppressed product is \eqref{fs:G9} with total count at most eight and target \(W_P^{1/2}\); we now put the high \(B^h\) in \(w_H^{1/2}\) and the lower one in point with \(Nw_H^{-1/2}\lesssim N r^{e/2}\tau^{-A_H/2}\). The \(\tau\) loss of \eqref{fs:G2} is covered. Coefficients through this total order cost at most \(C r^{-1+2\eta_*}\). The placements work also with plain \(O(r^{-1+3\eta_*})\) bounded coefficients, or with one \(B^h\) replaced by \(B^\psi\) if the total count on the two inputs is at most seven (measure a high \(B^\psi\) in its own level; point \(B^\psi\) through its stated orders in \eqref{fs:G1}). Thus \(T^jD\) on the vanishing-coefficient terms works simply counting \(D\) as up to one extra lab first derivative on the metric/field factors; differentiated \(s_s\) coefficients are bounded. Additional unscaled spatial lab derivatives can always be distributed before switching to \(Z,T\) counts, using \(\partial_d=\ell_d T+r^{-1}b_d Z\) on lab jets with \(T\) still commuting through other lab derivatives. Smooth bounded field compositions with explicit \(r^{-1}\) suppression and the analogous symbol bounds work by the same allocations.

Two explicit comparisons record the restrictive suppressed products.
At strong level a coefficient $O(r^{-1+3\eta_*})$,
a high $B^h$ measured with $w_H$, and a low $Nw_H^{-1/2}$ give
\[
 W_Pr^{-2+6\eta_*}w_H^{-2}
 \lesssim r^{p_P-2+6\eta_*+2e}\tau^{A-2A_H},
\]
whose radial exponent is $-.38+6.002\eta_*<0$.
At coefficient count nine only level $0$ occurs; its squared
comparison is bounded by
$r^{p_0+.021-.96+e}=r^{-.299+e}$.
For coefficient count at least seven the two input counts sum to
at most two, so both inputs admit two additional rotations.
For a lower coefficient the smaller input has count at most four
before those rotations.  The mixed strong product has total count
seven, and its lower input needs at most five counts after the
rotations, within the order-six point bound for $B^\psi$.

\par\medskip\noindent\textbf{Claim 4: the leading weak-null source and its corrections.}\label{fs:claim-quadratic}\quad
We treat \(Q_\eta\) separately. First take \(T^j\) on it, for either equation, giving the polarization sum. Set \(h_i=T^i h,\ p_i^\flat=T^{i+1}h\) (here this decorated letter denotes just the symmetric tensor, not a radial power), and set
\[
       F_i=r\Pi T^{i+1} h,\qquad \widehat F_i=r\Pi T^{i+1}\psi .
\]
For each polarization pair, use \eqref{fe:gauge-projection} to write
$\partial_\alpha h_i=\ell_\alpha\breve p_i+E_{\alpha i}$ with
$\mathcal G_n\breve p_i=0$ and $\Pi\breve p_i=\Pi p_i^\flat$.
The error \(r E_{\alpha i}\) is made of \(b Z h_i\) good terms and \(r\) times the differentiated gauge error \eqref{fe:F2} with bounded sphere coefficients. Here the lab differentiations in \eqref{fe:F2} are taken before the bad/flat-good decomposition. Expand each quadratic polarization sequentially, so each remainder need use just one such error and as its partner a full \(\partial h_l\) or a projected \(\ell\breve p_l\), estimated without any substitution of gauge errors into that partner. This gives
\begin{equation}
\label{fs:G10}
T^j Q_\eta
      =r^{-2}\ell\otimes\ell\sum_{i+l=j}\mathcal C_{il}(F_i,F_l)
                     +{\rm Err}_j,
\end{equation}
where the pairings include constants. One may use ambient spatial tensor components for the screen, the Euclidean pairing being smooth. After \(Z^{c_3}\) the \(r^2\)-scaled error terms have placements \eqref{fs:G9} or explicit \(O(r^{-1})\) suppressions of that type, \(\widetilde a\,B^h\), or
\begin{equation}
\label{fs:G11}
G^h_{b_1 i_1}\,B^h_{b_2 i_2},\qquad
                    b_1+i_1+b_2+i_2\le n_* .
\end{equation}
Smooth bounded compositions are allowed on suppressed terms. These assertions follow from \eqref{fe:F2}, including \(r\Theta=O(\widetilde a)\). Terms with \(\widetilde a\) fit directly by \eqref{fs:G2}. For \eqref{fs:G11} at levels \(H,0\), if the good input is lower order use its point \(C N r^{-1+2\eta_*}\tau^{-A_H/2}\) and measure the bad input in \(w_{i,i_2}\). If the good input is higher order measure it in \(W_i\), the low bad one in point \(C N w_H^{-1/2}\); this works in both weights since \(w_H^{-1}\lesssim1+M\), and \(w_{i,j}w_H^{-1}\lesssim W_i\).

In \(T^j D Q_\eta\) the commutator terms \((T^m s_s^d)\partial_d T^{j-m}Q_\eta,\ m\ge1\), already have the suppressed strong placement via \(T^m s_s^d=O(\widetilde a)\) with \(Z\)-bounds. Consider thus \(D\) on \eqref{fs:G10} before the \(Z^{c_3}\) distribution. When it differentiates \({\rm Err}_j\) the suppressed terms are still suppressed, now at total orders up to eight; expressions involving the gauge have the same differentiable product structure, with \(O(1/r)\) when a sphere projector is differentiated by \(D\). The good factor underlying \eqref{fs:G11}, if hit by \(D\), now costs \(C r^{-1} B^h_{b_1+1,i_1}\): indeed \(D=(1-s_s\cdot n)k^{-1}\partial_q+r^{-1}d Z\) with bounded coefficients under spatial-emitter differentiation, and \(\partial_q\) commutes with \(Z\) before expanding on \(h_i\). Differentiation of scale factors costs a relative \(O(1/r)\) too. Thus the remaining unsuppressed placements are \eqref{fs:G11} with total count up to eight (one extra derivative can now be on the bad input). If good is low the previous point placement and \(w_H\) on the bad input suffice even with \(W_P\) since \(2A_H>A\). If good is high, use \(W_H\) and the low bad point bound; the relevant comparisons are
\[
       (W_P/W_H)w_H^{-1}\lesssim1+M,\qquad
       (w_{P,j}/W_H)w_H^{-1}\lesssim1.
\]
The first uses \(p_P-p_H+e<.72\) within \(|q|\le r\); outside, use \(\mu\gtrsim r^{1-\delta}\tau\). The second follows from \eqref{fs:G2}. Low placements here have at most four derivatives before the two spare rotations. This proves acceptability of all the remainders of \eqref{fs:G10}.

Consider the outgoing budget of the first term of \eqref{fs:G10} for \(i_{\rm lev}=H,0\). On \(t\gg r\) ordinary high-low placement on differentiated \(F_i,F_l\) suffices with \(B^{i_{\rm lev}}\) for the high one, low \(w_H^{-1/2}\) in point for the other:
\[
 (W_{i_{\rm lev}}/w_{i_{\rm lev},i'})w_H^{-1}\lesssim1
                      \qquad(t\gg r,\ i'\ {\rm in\ that\ level}).
\]
On the remaining \(t\lesssim r\) use the mixed entry. A high differentiated \(F_i\) is controlled by the corresponding \(B_1\chi\)'s of \(h\) and the \(V^h_{b i}\) term in the cap comparison. The latter term can instead be charged as a high good input in bulk paired with low point \(N w_H^{-1/2}\), as in \eqref{fs:G11}. The \(B_1\)'s have time cap with square weight \(w_{i_{\rm lev},i}\). Thus the mixed kernel is bounded by
\[
      C N\sup(w_{i_{\rm lev},i}w_H)^{-1/2}
             \quad(\text{sup at fixed }\mathfrak u\text{ in the shell})
\]
on \(|u|\lesssim R_f\). Its squared integral for \(i_{\rm lev}=H\) costs \(C N^2 R_f^{2e+\delta}\). For \(i_{\rm lev}=0\) use \(w_{0,i}^{-1}\lesssim R_f^{.021}\), giving cost at most \(C N^2 R_f^{.021+e+\delta A_H+1-A_H}\). Here \(r^\delta+|u|\) is comparable to \(r^\delta+|q|\), and the low point placement including angular sup has the orders of \eqref{fs:G1}. Both costs give shell summability with the indicated \(p_{i_{\rm lev}}\).

The squared mixed-shell exponents just obtained are
\begin{equation}
\label{fs:quadratic-shell-margins}
 \begin{split}
 p_H-1+2e+\delta&=(-1/6+.012)\eta_*<0,\\
 p_0-1+.021+e+\delta A_H+1-A_H
   &=-.339+1.011\eta_*-.01\eta_*^2<0 .
 \end{split}
\end{equation}
Their strict signs justify summing the norm charges themselves
over dyadic shells.

In the strong outgoing budget for the first term of \(D\) applied to \eqref{fs:G10}, use the correction (one for each term with its actual \(\mathcal C_{il}\))
\begin{equation}
\label{fs:G12}
\mathcal P_*= Z^{c_3}\!\left[
           (1-s_s\cdot n)r^{-1}\ell\otimes\ell\,\mathcal C_{il}(F_i,F_l)
                              \right].
\end{equation}
Indeed \(kD=(1-s_s\cdot n)\partial_q+r^{-1}d' Z\), \(k\) here being the one in \eqref{fs:G8}; \(\partial_q\) commutes with \(Z\). The outgoing remainder after multiplying by \(r\) costs
\(C\widetilde a \sum|Z^{\le c_3}(F_i F_l)|+C r^{-1}\sum|Z^{\le c_3+1}(F_i F_l)|\)
with product derivatives allowed individually in the sums. This uses \(\partial_q(1-s_s\cdot n)=O(\widetilde a)\) with spatial-emitter derivatives. These are suppressed products through total order eight. In the bulk correction costs of \eqref{fe:F14}, all terms have at least \(r^{-1}\) in front of differentiated products of \(F_i,F_l\), up to the factor \(r d_g\lesssim r^{2\eta_*}\). Here \(Z^b F_i\) uses \(B^h_{b i}\), and \(\partial_q Z^b F_i\) uses \(B^h_{b,i+1}\) by commuting then expanding on \(T^{i+1}h\). Thus the total input count never exceeds eight in these \(B\) indices either; the \(r^{-1+2\eta_*}\) suppressed estimate with \(W_P^{1/2}\) proves sufficiency. At a time-slice cap the count in \eqref{fs:G12} itself is at most seven; make the same high-low comparisons with \(w_H\) on the high \(B_1\chi\)'s. A \(V^h\) term missing from their cap control is bounded by \eqref{fs:G1} directly; after the weights and low point factor this term has sphere \(L^2\) bounded by \(C N^2 r^{p_P/2-2+3\eta_*}\) with favorable \(\tau\) factor, summable in \(L^2(dr)\).

\par\medskip\noindent\textbf{Claim 5: raw leading products on the finite slab.}\label{fs:claim-raw}\quad
It remains to verify the raw time-multiplier budgets of the leading terms. At \(H,0\) they use, after \(r^2\) scaling, products \(F_i F_l\) with \(i+l=j\); at level \(P\), \(F_i\widehat F_l\) with \(i+l=j\), modulo suppressed terms, all still differentiated by \(Z\) to order \(c_3\). Indeed
\(D F_i=\widehat F_i+O(r^{-1})B^h_{0,i}+O(\widetilde a)\sum_{m\le i}B^h_{0,m}\)
as a differentiable product bound under \(Z\), by commuting \(D\) with lab times before expanding its spatial-frame coefficients. On \(|q|>r^\nu\), normal high-low placement works without using \eqref{fe:F11}: \(w_{i_{\rm lev},j}/w_{i_{\rm lev},j'}\lesssim l^{9d_*}\) in a level, and \(w_{P,j}/w_H\lesssim l^{A-A_H}\) for a high \(F\) in the strong product. The low \(w_H^{-1/2}\) point factors offset both ratios by \eqref{fs:G1}, \eqref{fs:G2}. High factors use their corresponding \(B\) in weighted bulk.

Outside the near band the squared coefficient needed for the
within-level transfer is at most $l^{9d_*}w_H^{-1}$; for a high
$h$ input in the strong product it is at most
$l^{A-A_H}w_H^{-1}$.  They are bounded because
\[
 \begin{split}
 e-(\nu-\delta)A_H&=-.089\eta_*+.09\eta_*^2<0,\\
 9d_*-A_H&=-1+1801\eta_*<0,\qquad
 A-2A_H=-.64+2\eta_*<0 .
 \end{split}
\]

On $|q|\le r^\nu$, a uniform bound on the undifferentiated news,
rather than radial losses in every low differentiated supremum, closes
the raw product estimate.  First replace $F_i$ by $T^iF_0$ and
$\widehat F_l$ by $T^l\widehat F_0$.  Since
$T^m(r\Pi)=O(1)$ for $m\ge1$, with $Z$ derivatives allowed, each
replacement error has $r^{-1}$ size relative to $B^h$ or $B^\psi$
at the allowed indices.  The remaining raw jet still contains a lab
$T$ derivative.  Products with these errors have the suppressed
estimates already proved, for both $h$ and $\psi$ inputs.

Tile the near region in $(q,\rho,n)$, with $\rho=\log r$, using
small fixed coordinate boxes and a fixed enlargement of bounded
overlap.  The measure identity is
\begin{equation}\label{fs:tile-measure}
 t=q+e^\rho,\qquad \frac{dt\,dr}{r}=dq\,d\rho.
\end{equation}
On the enlarged boxes, $|q|\lesssim r^\nu$, the operators $T,Z$
have uniformly bounded smooth coefficients through the required
orders, and the weights are comparable.  Within each level all time
counts have equivalent weights there; also $w_{P,j}\asymp w_H$.
Angular coordinates change the measure by a uniformly bounded factor.

Ordinary derivatives of $F_0$ through order $n_*$ are controlled in
the required weighted $L^2$ norm by $B^h_{bi}$ with $b+i\le n_*$;
the same statement holds for $\widehat F_0$ and $B^\psi$ through
the strong target order.  To verify this count, commute $\partial_q$
through $Z$, expand it on $T^{i+1}h$ successively as
$kT+r^{-1}s^db_dZ$, and retain the initial lab time.  Sphere rotations
span the angular derivatives.  High $h$ factors use their own level
in an $h$-only product and level $H$ in the strong product.

Write $\partial^a$ for the array of ordinary tile derivatives of
order $a$. For $a+b=m>0$, Hölder and Gagliardo--Nirenberg interpolation
on an enlarged tile give
\[
 \|\partial^af\,\partial^bg\|_2
 \le \|\partial^af\|_{2m/a}\|\partial^bg\|_{2m/b}
 \le CN\bigl(\|f\|_{H^m}+\|g\|_{H^m}\bigr),
 \qquad \|f\|_\infty+\|g\|_\infty\le CN.
\]
An order-zero factor uses the $L^\infty$ endpoint; the last inequality
is weighted arithmetic--geometric mean.  Insert cutoffs on the enlarged
tiles for inhomogeneous norms.  The needed supremum bound for
$f,g\in\{F_0,\widehat F_0\}$ is precisely \eqref{fe:F11}.
Local weight comparability, not a uniform bound on $\mu$ over the
whole near region, permits squaring and summing this estimate.

The same argument applies up to the current time endpoint.
For large $R_0$ the initial slice does not meet this band.  If the
upper endpoint meets a tile, it is the uniformly smooth graph
$\rho=\log(t_{\rm end}-q)$, with the domain below.  Extend across
it by a finite linear combination of values at negative multiples
of the signed graph-height, matching normal derivatives through
order $n_*-1$.  Change of variables gives uniform $H^{n_*}$ and
supremum bounds from enlarged interior boxes, and reflected weights
remain comparable.  Choose the fixed tile sizes small enough that,
for boxes meeting $r\ge R_0$, the required enlargements and reflected
points never extend inward beyond the fixed-factor overlap below
$R_0$; near-cutoff boxes may equivalently use compact overlap controls.
Source integrals with an additional
cap truncation can be bounded on the full time slab.  Summing with
bounded overlap proves the near-region raw budgets, with only the
stated compact-overlap and bootstrap-product costs.

\par\medskip\noindent\textbf{Claim 6: initial and starting-annulus costs.}\label{fs:claim-boundary}\quad
All corrections are added before forming \(\xi\) for each equation; the bounds concern \eqref{fe:F1} for uncorrected fields, already part of the total norm. On fixed starting overlap annuli their extra current and boundary costs satisfy the same compact convention: \eqref{fs:G5} corrections use a first derivative of a lower commuted field with bounded coefficients (weighted at the same time level, with acceleration decay where strong and weak inputs couple); in \eqref{fs:G7} smooth-acceleration corrections the coefficient carries \(\widetilde a\). The products in \eqref{fs:G12} on such annuli cost quadratically by weak fields with \(2A_H>A\). One can use the off-trap compact controls and sphere Sobolev placements there. Time-zero costs use the same cap estimates or directly the input symbol bounds. All patch/annulus constants on nonlinear terms are allowed fixed before smallness. This verifies \eqref{fs:F15} and therefore \eqref{fs:F}, with every positive-radius-gain term summable and the prior linear costs taken with independent far constants as specified there.

\end{proof}

\subsection{Simultaneous exterior closure}

The compact, far, and acceleration estimates concern the same finite
smooth slab.  Write $K=\mathsf K_R$, $F=\mathsf F$, and
$a_*=\mathsf S_a$, and collect the fixed data and superlinear
remainders in $E$.  The three estimates \Cref{co2:compact-high}, \eqref{fs:F},
and \Cref{co2:sliced-acceleration} have the form
\begin{equation}\label{fs:coupled-inequalities}
 \begin{split}
 K&\le\varepsilon_R(K+F)+E,\\
 F&\le\varepsilon_0(F+a_*)+B_0K+E,\\
 a_*&\le D_0K+C_0F+E .
 \end{split}
\end{equation}
Here $\varepsilon_0$ becomes small by increasing $R_0$, whereas
$\varepsilon_R$ is made small later by the compact matching radius
and plateau choices.  To choose them in this order, we need $C_0$
to stay bounded as $R_0$ grows.  The next lemma proves exactly that
fact and the compatible treatment of the pure acceleration defect.

\begin{lemma}[Radius-uniform coupling]
\label{fs:uniform-coupling}
The coefficient of $\mathsf F$ in \Cref{co2:sliced-acceleration} may be taken independent
of $R_0$.  The comparison of the background defect $f_B$ with $Q(0)a$
may be split at a radius $Y$ fixed before $R_0$: its compact part costs
$C(Y)\mathsf K_R$ and fixed data/nonlinear terms, and its tail costs
$C Y^{-1+\epsilon}\mathsf S_a$, with any sufficiently small fixed
$\epsilon>0$.  Conversely, the pure far defect, including the starting
annulus, has an $o_{R_0\to\infty}(1)$ coefficient on $\mathsf S_a$ in
\eqref{fs:F}.
\end{lemma}
\begin{proof}
On a fixed $t$ slice the emitter coordinates satisfy
\[
 y=C(t-r)-C(t)+rn,\qquad
 \partial_r y=n-s(q),\qquad \partial_Ay=r\partial_A n .
\]
It follows that
\begin{equation}
\label{fs:slice-jacobian}
\begin{gathered}
 d^3y=kr^2\,dr\,d\omega,\qquad
 \partial_{y_i}r=\frac{n_i}{k},\\
 \partial_{y_i}n_j=
 \frac1r\left(\delta_{ij}-\frac{(n_j-s_j)n_i}{k}\right),
 \qquad \operatorname{div}_y n=\frac2r .
\end{gathered}
\end{equation}
All comparison constants depend only on the fixed small speed bound.
Indeed
$(1-\|s\|_\infty)r\le r_y\le(1+\|s\|_\infty)r$.
These first derivatives contain no acceleration.  In particular one
must perform the single spatial integration by parts below in $d^3y$
before comparing it to the emitter measure.

Let $v_p$ be any adjoint zero mode in the fixed normalization of
\Cref{co2:sliced-acceleration}.  Its uniform symbol bounds are
$|v_p|\le C r_y^{-2+\epsilon}$ and
$|\partial_yv_p|\le C r_y^{-3+\epsilon}$.
The leading linear source pairing has the form
\[
 I(t)=\int \chi_t\chi_>v_{p(t)}
 \bigl[A(p(t))T\psi+B^d(p(t))\partial_{y_d}\psi\bigr]\,d^3y .
\]
The coefficient matrices and the laboratory/component transformation
are spatially constant on this slice.  The outer cutoff $\chi_>$ can
be chosen so its emitter support begins at $r\ge cR_0$; a fixed-factor
intermediate annulus is charged to $C(R_0)\mathsf K_R$.
Write $T\psi=L\psi-n^d\partial_{y_d}\psi$.  Then
\[
 I(t)=\int\chi_t\chi_>v_p A L\psi\,d^3y
 -\int\partial_{y_d}\bigl(\chi_t\chi_>v_p(B^d-An^d)\bigr)
                    \psi\,d^3y .
\]
Derivatives of either cutoff cost $C/r_y$ on its support.
Equation \eqref{fs:slice-jacobian} and the adjoint symbol bounds give
\[
 |I(t)|\le C\int_{r\ge cR_0}
 r^\epsilon\bigl(|rL\psi|+|\psi|\bigr)\,\frac{dr\,d\omega}{r}.
\]
Choose $\epsilon<\zeta<(p_P-A)/2=.13$.  Cauchy--Schwarz contributes
\[
 \left(\int_{cR_0}^\infty r^{-2(\zeta-\epsilon)}\frac{dr}{r}\right)^{1/2}
 =\frac{(cR_0)^{-(\zeta-\epsilon)}}{\sqrt{2(\zeta-\epsilon)}}.
\]
This is uniformly bounded as $R_0$ grows.  Since
$W_P=r^{p_P-A}(t+10r)^A$ and \eqref{fe:F1} controls both
$\psi=(r\psi)/r$ and $rL\psi=L(r\psi)-\psi$, we obtain
\[
 \|\langle t\rangle^{A/2}I\|_{L^2_t}\le C\mathsf F ,
\]
with $C$ independent of $R_0$.  Bounded initial times are included
in the fixed finite-time input.  No integration by parts in time,
nor any differentiation of $p(t)$, has been used.

Choose $Y$ beyond the fixed compact operator supports and the
transition where target smoothing reaches its prescribed far form.  The
compact defect comparison on $r_y\le2Y$ is given independently by
\Cref{co2:compact-defect}.  On its complement it is enough to bound the two terms
separately:
\[
 |f_B|+|Q(0)a|\le C r_y^{-2}
             \bigl(\widetilde a+|a(t)|\bigr).
\]
Use the majorant $A_\rho$ from \eqref{fs:acceleration-average}.
On the cutoff $\chi_t$, $t$ is a sufficiently large fixed multiple
of $\rho$, so \eqref{fs:average-bounds} gives
$\|\langle t\rangle^{A/2}A_\rho\mathbf1_{\{t\ge2C_2\rho\}}
\|_{L^2_t}\le C\mathsf S_a$, uniformly in $\rho$.
Pairing with $v_p$, converting $d^3y$, and
using Minkowski yields
\[
 C\mathsf S_a\int_Y^\infty \rho^{-2+\epsilon}\,d\rho
 \le C_\epsilon Y^{-1+\epsilon}\mathsf S_a .
\]
The omitted $(1-\chi_t)Q(0)a$ tail obeys the same bound directly.
Thus $Y$ is fixed before $R_0$ without producing a
$C(R_0)\mathsf S_a$ term.  Finally the direct shell estimates for
the pure defect in the proof of \Cref{fs:source-estimate} also apply
to the single shell $[R_0/2,R_0]$.  They prove the last assertion.
All convolution bounds remain valid up to an arbitrary finite
lifespan by extending $a$ by zero after the endpoint; causality
leaves the values inside the slab unchanged.
\end{proof}

\begin{corollary}[Global exterior estimates from the base neighborhood]
\label{fs:closure}
There is a neighborhood specified solely by sufficiently small
$p_{10}$ difference from the Kerr bridge such that the two reduced
computing-cylinder evolutions are global and smooth.  They obey
\Cref{co2:low-bound,co2:sliced-acceleration,co2:sliced-static,co2:compact-high} and \eqref{fs:F}, with arbitrarily
small total base bounds after the neighborhood is sufficiently
shrunk.  The higher smooth norms on growing cylinders have the
subexponential bounds of \Cref{co2:higher}; their constants may depend
on the fixed smooth datum and the order.  No smallness of those
higher initial norms is required.
\end{corollary}
\begin{proof}
Proposition~\ref{far:closure} proves the simultaneous absorption and
continuation from these estimates.  Its bootstrap cap can be chosen
arbitrarily small before the $p_{10}$ neighborhood is shrunk.  This
gives the asserted small base bounds on both cylinders.
The higher-order conclusion is then Proposition~\ref{co2:higher},
whose induction uses finite higher initial seminorms without shrinking
the common base neighborhood.
\end{proof}
